\markboth{DISCUSSION}{DISCUSSION}
\section{Discussion}\label{sec:discussion}

This section discusses architectural trade-offs, the scope boundaries we deliberately set, the limitations of the current system, and the practical lessons we drew from building Training Copilot.

\subsection{Architectural Trade-offs}\label{sec:discussion:tradeoffs}

\textbf{Multi-agent vs.\ single-agent design.} The multi-agent architecture introduces coordination overhead: each query traverses the orchestrator, one or more specialists, and their MCP servers, accumulating latency from multiple LLM calls and round-trips that a single-agent design would avoid. Three benefits justified this. With roughly 50 tools across all servers (before the optional Telegram bridge), a single agent frequently selected incorrect tools (research confirms larger tool sets increase misselection~\citep{qu_tool_2024, patil_gorilla_2023}), and scoping each specialist to at most three servers substantially improved accuracy. Each specialist also carries a domain-optimised prompt; a single prompt embedding all domain knowledge was unwieldy and often ignored, which the modular prompt architecture~\citep{wang_survey_2024} solved. Cross-domain queries also benefit from parallel execution, since concurrent \texttt{ask\_*} calls reduce latency to the slowest specialist rather than the sum.

\textbf{MCP vs.\ direct API integration.} Wrapping each API in an MCP server introduces a process boundary (every tool call is an HTTP round-trip) that direct integration would avoid, but the abstraction provides three properties that justify the cost, consistent with established microservice principles~\citep{dragoni_microservices_2017}: independence (each server manages its own authentication, rate limiting, and retry logic, so restarting one does not affect others), uniform discovery (adding the Telegram proxy and flythrough servers required one file and one configuration line each, confirming the single-line-extensibility promise), and deployment flexibility (servers are independently containerisable, with URLs swapped via environment variables for Docker Compose).

\subsection{Scope and Focus}\label{sec:discussion:scope}

Following the midterm review, we deliberately shifted from broadening the feature surface to deepening the capabilities at the heart of the product. Most subsequent effort therefore went not into new integrations but into hardening the two things we regard as Training Copilot's core: \textit{agentic data analytics} (turning raw, multi-source telemetry into synthesised, grounded insight) and \textit{agentic training and recovery planning}, in both senses of recovery: guidance for returning to training after injury, and day-to-day readiness inferred from physiological signals such as HRV, sleep, and Body Battery. This is where our evaluation puts its emphasis (Section~\ref{sec:evaluation}): the personas and scorers are built to stress the agentic reasoning path rather than the presentation layer. The remaining surfaces (the analytics dashboards, the Telegram delivery channel, the 3D flythrough, and the weather and calendar integrations) are intentionally lightweight add-ons, shaped to our own preferences as endurance athletes and easy to retune or replace; they demonstrate the architecture's extensibility rather than constitute the contribution, which is why their rough edges (Table~\ref{tab:known_issues}) are acceptable at this stage.

Training Copilot is a research prototype that tests two hypotheses: that MCP-based data integration achieves single-line extensibility, and that multi-agent coordination produces richer, more context-aware recommendations than a monolithic approach. Concretely, the path we prioritised is the \textit{core conversational analysis and planning workflow}, from a natural-language question through agent coordination to a data-grounded, synthesised answer, while accepting known incompleteness in peripheral features. The capabilities we consider production-quality for a research prototype are training-readiness assessment (combining recovery, weather, calendar, and activity data), training-load analysis (CTL/ATL/TSB trends), route planning with interactive maps, RAG-grounded fitness advice with source citations, schedule-aware planning with calendar event creation, and proactive, schedule-triggered goal check-ins. The rough edges we flagged at the midterm have since been addressed, and this pass surfaced no new ones; Table~\ref{tab:known_issues} is therefore a placeholder for a final pre-submission re-assessment of anything that remains.

\begin{table}[ht]
\begin{center}
\begin{minipage}{\textwidth}
\caption[Known incomplete areas]{Known incomplete areas (placeholder, pending a final pre-submission re-assessment).}\label{tab:known_issues}
\begin{tabularx}{\textwidth}{X}
\toprule
\textbf{Known incomplete areas} \\
\midrule
\textit{To be finalised before submission.} The rough edges noted at the midterm (inconsistent chart triggering, chat-map metric overlays, within-session context windowing, relative-date handling for calendar writes, and sport-aware route timing) have since been addressed, and no new ones surfaced in this pass. Any that remain at submission will be re-assessed and listed here. \\
\botrule
\end{tabularx}
\end{minipage}
\end{center}
\end{table}


\subsection{Limitations}\label{sec:discussion:limitations}

\textbf{Latency.} Complex cross-domain queries typically take 8-15 seconds end-to-end in our development setup, which is acceptable for training planning but not for in-workout guidance. LLM inference time dominates; MCP round-trips contribute under 500\,ms in aggregate. Token-level streaming would improve perceived responsiveness but is not yet enabled; the agents run non-streaming for robustness in the multi-agent loop.

\textbf{Knowledge currency.} The RAG-based fitness specialist draws from fifty public-domain books from Project Gutenberg, necessarily historical. While progressive overload, periodisation, and recovery principles are timeless~\citep{grgic_periodization_2017}, modern exercise science has advanced substantially, and contemporary open-access publications would improve the specialist's relevance. We also cannot yet confirm that retrieved passages are always relevant or faithfully used in the final answer; grounding is expected to reduce hallucination but is not a guarantee, which is why citation-faithfulness verification is planned (Section~\ref{sec:conclusion:next}).

\textbf{Platform dependencies.} Training Copilot depends on external APIs of varying availability. Strava now requires a subscription for Standard Tier API access~\citep{strava_api_update_2026}, and Garmin Connect has no official public API, so our server relies on a reverse-engineered third-party library that can break without warning when Garmin changes its internal endpoints; the MCP architecture mitigates this, since each server is independently replaceable, but does not eliminate the dependency.

\textbf{Proactivity depends on one process.} The scheduler that fires check-ins, daily nudges, and deep-analysis reports lives in the always-on Telegram bridge (the only long-running process), so if that process is down, proactive delivery stops silently until it restarts. Moving the scheduler into a dedicated service would remove this single point of failure; the code is structured to allow it later.

\textbf{Route quality and data availability.} Some of the most popular consumer route planners, such as Komoot and Outdooractive, keep their curated, community-contributed route catalogues behind closed doors: that route data \textit{is} their product, so declining to expose it through a public API is a rational commercial choice rather than an oversight. Training Copilot therefore plans over OpenRouteService, an open engine on raw OpenStreetMap data; it can synthesise a loop of any target distance from any start point, even where no pre-planned route exists (Section~\ref{sec:data:other}), but it optimises over the road and path graph rather than serving hand-picked scenic itineraries, so a suggestion's scenic quality can trail a Komoot route. This is a limitation of the available data ecosystem, not of our integration, which would accept a richer route provider through the same one-server interface the moment one opens up. Because the project's focus is agentic analytics and training/recovery planning, the search for a higher-quality route source is paused for now and left to future work.

\textbf{Single-user focus.} Multi-tenant deployment with per-user credential vaults and session isolation is architecturally supported but not yet implemented; credentials currently sit in a plaintext \texttt{.tokens/} directory. The security implications (SSRF prevention, tool-output sanitisation, egress control for user-added MCP servers, encrypted token storage) remain open; the MCP-specific threat surface that \cite{hou_mcp_2025} catalogue (tool poisoning, prompt injection, and lifecycle attacks on third-party servers) applies directly once users can register their own servers, and it must be handled before any public deployment.

\subsection{Ethical Considerations}\label{sec:discussion:ethics}

An AI system that provides training recommendations bears responsibility for athlete safety. Training Copilot explicitly does not provide medical advice; system prompts instruct all specialists to recommend professional consultation for injury-related questions or abnormal metrics. Consumer wearables have known accuracy limitations~\citep{fuller_reliability_2020}, mitigated partially through multi-source corroboration (cross-referencing HRV, sleep, and Body Battery rather than any single metric) but not eliminated. The fitness RAG corpus also reflects early 20th-century biases and limited demographic scope; expanding it with contemporary, demographically inclusive literature is a priority for future work.

\subsection{Lessons Learned}\label{sec:discussion:lessons}

\textbf{Tool docstrings are the interface.} A tool's docstring quality has an outsized effect on selection accuracy: vague descriptions led to incorrect calls, while precise descriptions that state \textit{when} to call a tool substantially improved accuracy, consistent with \cite{patil_gorilla_2023}'s findings on API documentation quality. We treated docstrings as the primary interface between model and data layer.

\textbf{Graceful degradation is not optional.} With eight external dependencies, partial outages are the norm. The three-layer strategy (skip missing MCP servers, report unavailable specialists, surface tool errors as structured data) proved essential when individual services restarted or returned partial data during development; what we first treated as a nice-to-have quickly became load-bearing.

\textbf{Record full, clip for the model.} The recording-and-clipping pattern (Section~\ref{sec:agents:recording}) was one of our most impactful design decisions. GPS streams and intraday timelines can exceed 100\,KB; rendering the UI from the full trace rather than the model's response preserves data fidelity without burdening the LLM. We recommend this pattern for any tool-augmented agent that handles large structured results.

\textbf{Scoping constrains and enables.} Restricting each specialist to its own MCP servers began as a simplification but proved a design strength: it prevents tool confusion, reduces prompt complexity, and makes each specialist independently testable.

\textbf{Source preservation requires explicit engineering.} The orchestrator's synthesis step dropped citations in roughly one-third of fitness-related answers until we added an explicit \texttt{ensure\_sources()} post-processing step. Automated synthesis and citation preservation pull against each other, and closing that gap takes deliberate engineering.
